% VI23_FULL_SOLUTION_REFINEMENT_V1
\subsection*{VI.23.P13: The graph in affine coordinates}
\begin{problem}\label{prob:vi23-p13}
Let $f:\Spec B\to\Spec C$ be a morphism over $\Spec A$, corresponding to $C\to B$.
Identify the graph map on rings.
\end{problem}

\ifdefined\FullProblemDossiers
\paragraph{Why this problem matters.}
The graph formula is a useful bridge between the abstract pullback description and explicit coordinate algebra.

\paragraph{Useful ideas before starting.}
Translate the requested statement into compatible maps from an arbitrary test scheme, and use uniqueness before choosing coordinates.

\paragraph{Guided hints.}
\begin{enumerate}
\item Identify the product ring.
\item Record what the graph does after each projection.
\item Use the tensor-product universal property.
\end{enumerate}

\begin{solution}
Write $\alpha:C\to B$ for the $A$-algebra map defining $f$.  The target of the graph is
\[
\Spec B\times_{\Spec A}\Spec C=\Spec(B\otimes_AC).
\]
Since $\Gamma_f$ has components $\id_{\Spec B}$ and $f$, its ring map
$\gamma:B\otimes_AC\to B$ must restrict to the identity on the first factor and to $\alpha$ on
the second.  The tensor-product universal property therefore gives
\[
\gamma(b\otimes c)=b\,\alpha(c).
\]
Indeed
\[
\gamma(b\otimes1)=b,\qquad \gamma(1\otimes c)=\alpha(c).
\]
These are exactly the contravariant equations saying that the two graph components are
$\id$ and $f$.  Uniqueness of $\gamma$ is equivalent to uniqueness of the graph map into the
fiber product.
\end{solution}

\paragraph{What happened mathematically.}
A geometric assertion was reduced to a representability statement, so the canonical map was forced by universality.

\paragraph{Extensions.}
Repeat the argument after restricting to affine opens and compare with the tensor-product formula.

\paragraph{Provenance.}
\textbf{Canonical graph computation.}
\else
\paragraph{Provenance.} \textbf{Canonical graph computation.}
\fi
